\documentclass[11pt]{article}
\usepackage[T1]{fontenc}
\usepackage{lmodern,amsmath,amssymb,amsthm,geometry,enumitem,graphicx,xcolor,tikz,hyperref,fancyhdr}
\geometry{margin=0.9in}
\definecolor{courseblue}{HTML}{285E80}
\hypersetup{colorlinks=true,linkcolor=courseblue,urlcolor=courseblue}
\pagestyle{fancy}\fancyhf{}\fancyhead[L]{MATH\& 142 -- Precalculus II}\fancyhead[R]{Dr. Lee}\fancyfoot[C]{\thepage}
\newtheoremstyle{course}{7pt}{7pt}{}{}{\bfseries\color{courseblue}}{.}{.5em}{}
\theoremstyle{course}\newtheorem{definition}{Definition}[section]\newtheorem{example}{Example}[section]
\newenvironment{practice}{\par\medskip\noindent\textbf{Try it. }}{\par\medskip}
\setlist[enumerate]{itemsep=5pt}

\newcommand{\work}{\vspace{1.15in}}
\begin{document}
\begin{center}{\LARGE MATH\& 142 --- Midterm Examination}\\[5pt]
{\large Modules 1--5: Trigonometry and Modeling}\end{center}
Name:\ \rule{3in}{.4pt}\hfill Date:\ \rule{1.4in}{.4pt}

\smallskip\noindent\textbf{Instructions.} Show work and explain your reasoning. Give exact answers unless the question requests a decimal approximation. Angles are in radians unless marked with a degree symbol. A non-CAS scientific calculator is permitted. Total: 100 points.
\begin{enumerate}
\item \textbf{Angles and arc length (10 points).}
Convert $225^\circ$ to radians. Then calculate the arc length traced by a point on a wheel of radius $6$ cm through an angle of $5\pi/6$ radians.\work
\item \textbf{Unit circle (10 points).}
Determine exactly $\sin(7\pi/6)$, $\cos(7\pi/6)$, and $\tan(7\pi/6)$. State the reference angle and justify the signs.\work
\item \textbf{Six trigonometric functions (10 points).}
Suppose $\sin\theta=3/5$ and $\theta$ lies in Quadrant II. Find $\cos\theta$, $\sec\theta$, and $\cot\theta$ exactly.\work
\item \textbf{Graph transformations (10 points).}
For $f(x)=-3\cos(2(x-\pi/4))+1$, determine amplitude, period, phase shift, midline, and range. Explain the meaning of the negative amplitude coefficient.\work
\item \textbf{Modeling a rotating wheel (10 points).}
A rider starts at the rightmost point of a counterclockwise wheel of radius $9$ m, centered $11$ m above ground. The wheel completes a revolution every $36$ seconds. Write a model for the rider's height $h(t)$, and calculate the height at $t=9$ seconds.\work
\item \textbf{Inverse trigonometry (10 points).}
Evaluate exactly $\arcsin(\sin(5\pi/6))$. Then find $\cos(\arcsin(3/5))$ exactly and explain the principal-value restriction.\work
\item \textbf{Trigonometric identity (10 points).}
Prove $\dfrac{1-\cos^2x}{\sin x}=\sin x$ and state the set of inputs for which the displayed equality is defined.\work
\item \textbf{Solving an equation (10 points).}
Find \emph{all} solutions of $\sin(2x)=\sqrt{3}/2$ on $[0,2\pi)$. Show how you account for multiple cycles.\work
\item \textbf{Law of Sines (10 points).}
In triangle $ABC$, $A=40^\circ$, $B=65^\circ$, and $a=12$ cm. Determine $C$ and side $b$. Round the side length to two decimal places.\work
\item \textbf{Law of Cosines (10 points).}
Two sides of a triangle are $8$ m and $11$ m, with included angle $60^\circ$. Find the third side exactly and approximately to two decimal places.\work
\end{enumerate}
\end{document}
